\documentclass[11pt]{article}
\usepackage[margin=1in]{geometry}
\usepackage{amsmath,amssymb}
\newcommand{\finalanswer}[1]{\par\medskip\noindent\fbox{\parbox{0.94\linewidth}{\textbf{Answer:} #1}}}
\title{STAT 412 Homework 3 --- Question 11}
\author{}
\date{}
\begin{document}
\maketitle

\section*{Problem}
Impurities in batches of product of a chemical process reflect a serious
problem. It is known that the proportion of impurities $Y$ in a batch has the
following density function. Complete parts (a) through (e) below.
\[
f(y)=
\begin{cases}
12(1-y)^{11},&0\le y\le1,\\
0,&\text{elsewhere}.
\end{cases}
\]

(a) Verify that the given function is a valid density function. Select the
correct choice below and, if necessary, fill in the answer boxes within your
choice.

\begin{itemize}
\item[A.] The given function is a valid density function because it is at least
0 everywhere and because
\[
\int_0^1 12(1-y)^{11}\,dy=\fbox{\strut\hspace{1.5em}}.
\]
\item[B.] The given function is a valid density function because it is at least
0 everywhere.
\item[C.] The given function is a valid density function because
\[
\int_0^1 12(1-y)^{11}\,dy=\fbox{\strut\hspace{1.5em}}.
\]
\end{itemize}

(b) A batch is considered not acceptable if its percentage of impurities
exceeds 55\%. Find $P(Y>0.55)$. Round to four decimal places.

(c) What are the parameters $\alpha$ and $\beta$ of the beta distribution
illustrated here? \textit{(Type an integer or a decimal. Do not round.)}

(d) The mean of the beta distribution is $\dfrac{\alpha}{\alpha+\beta}$. What is
the mean proportion of impurities in the batch?
\textit{(Round to four decimal places.)}

(e) The variance of a beta-distributed random variable is
$\sigma^2=\dfrac{\alpha\beta}{(\alpha+\beta)^2(\alpha+\beta+1)}$. What is the
variance of $Y$ in this problem? \textit{(Round to six decimal places.)}

\section*{Work}
The function is nonnegative on $[0,1]$, and
\[
\int_{-\infty}^{\infty}f(y)\,dy
=\int_0^1 12(1-y)^{11}\,dy
=\left[-(1-y)^{12}\right]_0^1=1.
\]
Therefore $f$ is a valid probability density.

For part (b),
\[
P(Y>0.55)=\int_{0.55}^1 12(1-y)^{11}\,dy
=(1-0.55)^{12}=0.45^{12}\approx0.00006895\approx0.0001.
\]

For part (c), the beta density is
$\dfrac{1}{B(\alpha,\beta)}\,y^{\alpha-1}(1-y)^{\beta-1}$. Matching
$f(y)=12\,y^{0}(1-y)^{11}$ gives $\alpha-1=0$ and $\beta-1=11$, i.e.\
$\alpha=1$, $\beta=12$. (Check: $\tfrac{1}{B(1,12)}=\tfrac{\Gamma(13)}{\Gamma(1)\Gamma(12)}
=\tfrac{12!}{11!}=12$, the leading constant.)

For part (d),
\[
\mu=\frac{\alpha}{\alpha+\beta}=\frac{1}{1+12}=\frac{1}{13}\approx0.0769.
\]

For part (e),
\[
\sigma^2=\frac{\alpha\beta}{(\alpha+\beta)^2(\alpha+\beta+1)}
=\frac{1\cdot12}{13^2\cdot14}=\frac{12}{2366}\approx0.005072.
\]

\finalanswer{(a) Select Choice A and enter $1$.
(b) $P(Y>0.55)=0.0001$.
(c) $\alpha=1$, $\beta=12$.
(d) $\mu=1/13\approx0.0769$.
(e) $\sigma^2=12/2366\approx0.005072$.}
\end{document}
